\RevisionAddEnd
\begin{figure*}[t]
\centering
\resizebox{\textwidth}{!}{
\begin{tikzpicture}[
    node distance=1.15cm and 1.35cm,
    every node/.style={font=\small},
    process/.style={rectangle, draw, rounded corners, align=center,
      minimum width=3.15cm, minimum height=0.82cm, fill=blue!4},
    decision/.style={diamond, draw, aspect=2.2, align=center,
      inner sep=1pt, fill=orange!8},
    data/.style={rectangle, draw, double, align=center,
      minimum width=3.0cm, minimum height=0.82cm, fill=green!6},
    arrow/.style={-{Latex[length=2.3mm]}, thick}
]
\node[data] (input) {Binary cover $I$};
\node[process, right=of input] (partition) {Complete $3\times3$ blocks\\raster order; preserve borders};
\node[process, right=of partition] (hist) {Histogram $H$\\PEAKs $H(d)\ge2$; absent ZEROs};
\node[process, below=of hist] (rank) {For each PEAK, rank unassigned\\distance-one ZEROs with the CNN};
\node[process, left=of rank] (inject) {Assign one ZERO per PEAK\\and enforce global disjointness};
\node[process, left=of inject] (net) {Sort PEAKs by frequency\\select prefix maximizing\\$C_{\rm net}^{\rm map}$};
\node[data, below=of net] (message) {Payload $M$ and declared length $q$};
\node[process, right=of message] (embed) {Raster scan: $0\mapsto P_i$\\$1\mapsto Z_i$};
\node[data, right=of embed] (stego) {Stego image $I'$};
\node[data, below=of inject, yshift=-2.0cm] (wire) {Enumerative auxiliary stream\\$A_{\rm wire}=(q,r(\mathcal P),s_1,\ldots,s_t)$};

\draw[arrow] (input) -- (partition);
\draw[arrow] (partition) -- (hist);
\draw[arrow] (hist) -- (rank);
\draw[arrow] (rank) -- (inject);
\draw[arrow] (inject) -- (net);
\draw[arrow] (net) -- (message);
\draw[arrow] (message) -- (embed);
\draw[arrow] (embed) -- (stego);
\draw[arrow] (embed.south) -- (wire.north);
\end{tikzpicture}
}
\caption{Evaluated embedding pipeline. Learned distance-one assignment is
followed by explicit net-capacity optimization and enumerative serialization
within a reversible-data-hiding synchronization model.}
\label{fig:embedding-flow}
\end{figure*}

\begin{algorithm*}[t]
\caption{Net-Aware CNN-Ranked Distance-One Embedding}
\label{alg:embedding}
\begin{algorithmic}[1]
\Require Binary image $I$, payload $M=(m_1,\ldots,m_q)$, trained ranker $f_\theta$
\Ensure Stego image $I'$, serialized auxiliary stream $A_{\rm wire}$
\State Enumerate complete $3\times3$ blocks and compute $H(d)$
\State $\mathcal P_0\gets\{d:H(d)\ge2,\ d\notin\{0,511\}\}$; sort by $(-H(d),d)$
\State $\mathcal Z\gets\{d:H(d)=0\}$; $\mathcal U\gets\emptyset$; $\mathcal T\gets\emptyset$
\For{each $P_i\in\mathcal P_0$}
    \State $\mathcal C_i\gets\{Z\in\mathcal Z\setminus\mathcal U:D(P_i,Z)=1\}$
    \If{$\mathcal C_i\neq\emptyset$}
        \State $Z_i\gets\arg\min_{Z\in\mathcal C_i}(f_\theta(I,P_i,Z),Z)$
        \State $\mathcal T[P_i]\gets\{0\mapsto P_i,1\mapsto Z_i\}$;
        $\mathcal U\gets\mathcal U\cup\{Z_i\}$
    \EndIf
\EndFor
\State Select the frequency-ordered prefix $\mathcal T^\star$ maximizing
\RevisionAddBegin
$C_{\rm net}^{\rm map}(\mathcal S)=
\sum_{P_i\in\mathcal S}H(P_i)-|A_{\rm wire}(\mathcal S)|_2$
\RevisionAddEnd
\If{$q>\sum_{P_i\in\mathcal T^\star}H(P_i)$}
    \State reject the requested payload
\EndIf
\State $I'\gets I$; $j\gets1$
\For{each complete block of $I$ in raster order}
    \If{its pattern is $P_i\in\mathcal T^\star$ and $j\le q$}
        \State replace it by $\mathcal T^\star[P_i][m_j]$; $j\gets j+1$
    \EndIf
\EndFor
\State Serialize $q$, the active PEAK subset rank, and base-9 flip positions
\State \Return $I'$, $A_{\rm wire}$
\end{algorithmic}
\end{algorithm*}

\subsection{Efficient Extraction and Restoration}
\RevisionAddBegin
Extraction is designed as the exact inverse of the embedding contract, using
the marked image and the serialized auxiliary stream. The auxiliary stream
is decoded first because it fixes three pieces of synchronization information:
the payload length, the active PEAK subset, and the single flip position that
reconstructs each assigned ZERO. The decoder uses the serialized tables
directly, making the transmitted synchronization contract the source of
recovery decisions.

Once the active tables are reconstructed, they are collapsed into a single
inverse dictionary. This is possible because the encoder enforces global
injectivity: every stego pattern used by the codec belongs to at most one active
mapping. Therefore, observing a PEAK pattern decodes bit 0 and restores the same
PEAK, while observing its assigned ZERO decodes bit 1 and is also restored to
the PEAK. Patterns outside the dictionary are copied unchanged. This dictionary
view turns extraction into one linear raster scan with constant-time lookup.

The serialized payload length provides the stopping condition. After $q$ bits
have been recovered, remaining blocks are left unchanged because embedding used
the same declared stopping point. If the scan finishes before $q$ bits
are recovered, the pair $(I',A_{\rm wire})$ is rejected as inconsistent. This
explicit rejection rule is necessary because exact reversibility is guaranteed
for matched marked-image and auxiliary-stream pairs transmitted over a lossless
channel.

\RevisionAddEnd
\begin{figure*}[t]
\centering
\resizebox{\textwidth}{!}{
\begin{tikzpicture}[
    node distance=1.15cm and 1.35cm,
    every node/.style={font=\small},
    process/.style={rectangle, draw, rounded corners, align=center,
      minimum width=3.15cm, minimum height=0.82cm, fill=blue!4},
    decision/.style={diamond, draw, aspect=2.2, align=center,
      inner sep=1pt, fill=orange!8},
    data/.style={rectangle, draw, double, align=center,
      minimum width=3.0cm, minimum height=0.82cm, fill=green!6},
    arrow/.style={-{Latex[length=2.3mm]}, thick}
]
\node[data] (stego) {Stego image $I'$};
\node[data, below=of stego] (wire) {Auxiliary stream $A_{\rm wire}$};
\node[process, right=of wire] (decode) {Decode $q$, PEAK subset rank,\\and base-9 flip positions};
\node[process, above=of decode] (inverse) {Build globally injective inverse\\$P_i\mapsto(0,P_i)$, $Z_i\mapsto(1,P_i)$};
\node[process, right=of inverse] (scan) {Scan complete $3\times3$ blocks\\in raster order};
\node[decision, below=of scan] (done) {$q$ bits\\recovered?};
\node[process, right=of done] (lookup) {If pattern is in inverse:\\append bit and restore $P_i$};
\node[data, below=of lookup] (outputs) {Recovered payload $M$\\and exact cover $I$};
\draw[arrow] (stego) -- (inverse);
\draw[arrow] (wire) -- (decode);
\draw[arrow] (decode) -- (inverse);
\draw[arrow] (inverse) -- (scan);
\draw[arrow] (scan) -- (done);
\draw[arrow] (done) -- node[above]{Continue} (lookup);
\draw[arrow] (lookup) |- (scan);
\draw[arrow] (done.south) |- node[pos=0.25,left]{Yes} (outputs.west);
\end{tikzpicture}
}
\caption{Extraction and restoration pipeline. The serialized payload length
provides an unambiguous stopping condition, and incomplete border blocks remain
untouched.}
\label{fig:extraction-flow}
\end{figure*}

To reduce extraction complexity, all mapping tables are merged into a single inverse dictionary:
\begin{equation}
\mathcal{D} : p \mapsto (m, P_i),
\label{eq:inverse-dictionary}
\end{equation}
where $p$ is a stego pattern, $m\in\{0,1\}$ the extracted bit, and $P_i$ the
original pattern.

\begin{algorithm*}[t]
\caption{Deterministic Extraction and Exact Restoration}
\label{alg:extraction}
\begin{algorithmic}[1]
\Require Stego image $I'$, serialized auxiliary stream $A_{\rm wire}$
\Ensure Extracted message $M$, restored image $I$
\State Decode payload length $q$, active PEAKs, and their flip positions
\State Reconstruct all tables $T_i=\{0\mapsto P_i,1\mapsto Z_i\}$
\State Build $\mathcal D[P_i]=(0,P_i)$ and $\mathcal D[Z_i]=(1,P_i)$
\State $I\gets I'$; $M\gets\emptyset$
\For{each complete block of $I'$ in raster order}
    \If{$|M|=q$}
        \State \textbf{break}
    \EndIf
    \State Compute pattern index $d$
    \If{$d \in \mathcal{D}$}
        \State Append the decoded bit to $M$ and replace the block by its PEAK
    \EndIf
\EndFor
\If{$|M|\neq q$}
    \State reject the inconsistent stego/auxiliary pair
\EndIf
\State \Return $M$, $I$
\end{algorithmic}
\end{algorithm*}

\subsection{Learned Ranking, Ablations, and Wire Accounting}

\subsubsection*{Distortion characteristics}

For each embedded block mapped from $P_i$ to $\mathcal{Z}_i$, the number of flipped pixels is bounded by:
\begin{equation}
\max_{Z_j \in \mathcal{Z}_i} D_{i,j},
\label{eq:distortion-bound}
\end{equation}
ensuring pattern-specific distortion control.

\subsubsection*{Gaussian-process diagnostic}
\label{subsec:param_optim}

The parameter $\alpha$ is explored by Gaussian-process Bayesian optimization.
For every evaluated value, the implementation records capacity, PSNR, DRD, and
the normalized objective
\begin{equation}
J(\alpha)=\lambda\left(1-\frac{C_{\text{total}}(\alpha)}{C_{\mathrm{ref}}}\right)
+(1-\lambda)\mathrm{DRD}(\alpha).
\label{eq:gp-objective}
\end{equation}
Expected Improvement selects the next value of $\alpha$
\cite{Frazier2018Bayesian}. The equal weighting
$\lambda=0.5$ gives the two normalized diagnostic terms identical influence;
\RevisionAddBegin
it serves only the exploratory scalarization of the unconstrained diagnostic.
A $\lambda$ sweep therefore characterizes the exploration stage, whereas all
final tables use the independently specified distance-one policy. Optimizing
$\alpha$ alone provides incomplete control of
\RevisionAddEnd
$\mathrm{DRD}<1$, motivating the constrained distance-one policy used by the
\RevisionAddBegin
integrated system. The adaptive-threshold parameter $\alpha$ belongs to this
unconstrained diagnostic; the deployed codec fixes the policy to Hamming
distance one, and the final operating tables are governed by that fixed policy.
The exploratory search uses
\RevisionAddEnd
$\alpha\in[0,1.5]$, initial points $\{0,0.5,1,1.5\}$, six EI iterations on a
\RevisionAddBegin
151-point grid, $\lambda=0.5$, exploration 0.01, and a Mat\'ern-$5/2$ kernel
\RevisionAddEnd
with a white-noise term. The random seed is 20260607.

\subsubsection*{CNN-based candidate ranking}

For a PEAK pattern and each distance-one ZERO candidate, the CNN receives three
\RevisionAddBegin
$9\times9$ channels: the centered PEAK, the centered candidate, and the average
\RevisionAddEnd
spatial context of all PEAK occurrences. It regresses the empirical
reciprocal-distance replacement cost. The candidate with the smallest predicted
\RevisionAddBegin
cost is assigned to symbol 1, while symbol 0 preserves the PEAK.
\RevisionAddEnd

The $9\times9$ support covers the central $3\times3$ block and a three-pixel
margin in every direction. It contains the complete $5\times5$ DRD support
around each of the nine possible flipped pixels, including pixels outside the
\RevisionAddBegin
embedding block; a $7\times7$ patch covers corner flips only partially.
Larger patches add context beyond the local training target.
\RevisionAddEnd

The network contains two $3\times3$ convolutional layers with 16 and 32
channels and ReLU activations, adaptive $3\times3$ average pooling, and a
$288\rightarrow64\rightarrow1$ regressor with ReLU and Softplus. Training uses
Smooth-L1 loss, AdamW with learning rate $10^{-3}$, batch size 128, 20 epochs,
\RevisionAddBegin
and seed 20260607. Complete source images define the train/validation split at
image level.
\RevisionAddEnd

Figure~\ref{fig:cnn-architecture} summarizes the evaluated architecture and
its scalar candidate-cost output.

\begin{figure*}[t]
\centering
\resizebox{\textwidth}{!}{
\begin{tikzpicture}[
  node distance=0.7cm,
  layer/.style={rectangle,draw,rounded corners,align=center,
    minimum height=0.75cm,minimum width=2.3cm,fill=blue!5},
  arrow/.style={-{Latex[length=2mm]},thick}
]
\node[layer] (input) {Three $9\times9$ channels\\PEAK, ZERO, mean context};
\node[layer, right=of input] (conv1) {$3\times3$ Conv\\3$\rightarrow$16, ReLU};
\node[layer, right=of conv1] (conv2) {$3\times3$ Conv\\16$\rightarrow$32, ReLU};
\node[layer, right=of conv2] (pool) {Adaptive average pool\\$32\times3\times3$};
\node[layer, right=of pool] (reg) {Flatten, FC 288$\rightarrow$64\\ReLU, FC 64$\rightarrow$1, Softplus};
\draw[arrow] (input) -- (conv1);
\draw[arrow] (conv1) -- (conv2);
\draw[arrow] (conv2) -- (pool);
\draw[arrow] (pool) -- (reg);
\end{tikzpicture}}
\caption{Candidate-cost CNN used to rank globally available distance-one ZERO
patterns. The scalar output estimates reciprocal-distance replacement cost.}
\label{fig:cnn-architecture}
\end{figure*}

\subsubsection*{Random-Forest uniform-block ablation}
\label{subsec:rf-layer}

Uniform blocks are processed by a second reversible layer. A Random Forest
classifier predicts whether a block admits a low-cost single-pixel flip from
its $9\times9$ context, and a second forest predicts the pixel position. The
\RevisionAddBegin
safety label uses a local reciprocal-distance cost at most 0.5.
\RevisionAddEnd
Selected coordinates and positions are compactly serialized for exact
extraction and restoration. The experiment shows that the classifier
identifies reliable locations, while the coordinate stream defines the next
compression target. The final net-aware configuration therefore assigns its
bit budget to the enumerative block-mapping layer.

